


\section{Initial conditions}{\label{sect_in}}

We denote the initial moment in the laboratory frame by \(t_0\) and the initial position of the particle by \({\pmb{\mathfrak R}}_0\). We assume that at \(t_0\) the particle velocity is \({\bf v}_0\);  the corresponding four-velocity will be denoted by \( u_0\), and the four-momentum by \( p_0\). The field propagates along the \({\bf n}\) direction, chosen parallel to \(Oz\), and is described by the potential \( A(\phi)\), where \(\phi = ct - z\). We denote by  \({ A}_0\)  the value of the  potential \({ A}\) in the initial position of the particle \({pmb{\mathfrak R}_0}\)  at the initial moment \(t_0\) i.e.
\begin{equation}
    A_0\equiv(0,{\bf A_0})= A(\phi_0), \qquad \phi_0=ct_0 - {\pmb{\mathfrak R}_0}\cdot {\bf n}.
\end{equation}
For the case of a pulse, the most natural choice for \(t_0\) is a moment sufficiently far in the past, when the pulse has not reached the origin  and the particle is free (i.e. \( A_0= 0\)).
The general equations of motion will be written for arbitrary values of \(t_0\) and \({\bf v}_0\),particular results will be presented for two cases of interest:


\begin{enumerate}
    \item Gaussian pulse; the particle is at rest in \({\pmb{\mathfrak R}}_0\) at a moment \(t_0\) very far in the past, when the pulse has not yet reached the point \({\pmb{\mathfrak R}}_0\).
    \item monochromatic field; at the moment \(t_0\)  the particle is located in \({\pmb{\mathfrak R}}_0\), with an initial velocity \({\bf v}_0\), the value of the vector potential in the initial position of the particle being \(A_0= A(ct_0 - {\pmb{\mathfrak R}}_0 \cdot {\bf n})\).
    \item We will combine the two cases into a general one, corresponding to a field with a gaussian turn-on, followed by a monochromatic oscillation of infinite length. The initial conditions are the same as in the previous case (1).
\end{enumerate}

